% =============================================================================
% 7 CONCLUSÃO — primeira versão (18/11/2026), escrita com a skill
% revisao-conclusoes-cientificas: só conclusões já demonstradas na
% Discussão, no presente, sem "tendências" (gpt-oss-120b: sem diferença
% significativa após Holm), sem especulação.
% Guia UECE p. 66: seção numerada.
% Fontes: monografia/textuais/06-discussao.tex; docs/briefing.md §7;
% docs/error_taxonomy.md (implicações).
% =============================================================================
\chapter{Conclusão}
\label{cap:conclusao}

Em agentes que geram SQL livremente a partir de uma descrição textual do
banco, o que aumenta a correção das consultas são as regras sobre como a
base deve ser consultada, e não as descrições do que cada coluna
significa. Descrições de catálogo resolvem a escolha entre colunas
parecidas, mas não alteram de forma significativa a acurácia do agente.
Regras de relacionamento e de formato da resposta resolvem os
erros mais frequentes, que produzem consultas válidas e respostas
erradas, e elevam a acurácia de forma significativa no
\arq{qwen3.8-27b}. No \arq{gpt-oss-120b}, as três condições não diferem
estatisticamente com o banco de 40 perguntas.

As regras que mais contam são as que tratam das imperfeições dos dados e
da forma da resposta: o tratamento de referências órfãs e a restrição às
colunas pedidas. Essas informações não pertencem a nenhuma coluna
específica e, por isso, não cabem em um catálogo de dados. Nas camadas
semânticas industriais, elas são garantidas pelo compilador que gera o
SQL; em um agente que escreve o SQL, precisam ser dadas como instrução.

A forma mais simples de camada semântica, um arquivo de texto anexado ao
prompt, basta para reproduzir o padrão relatado pela indústria com
plataformas proprietárias: o catálogo acrescenta pouco, e o maior ganho
vem da camada semântica. Avaliar e adotar esse tipo de contexto não exige
licença paga, catálogo proprietário nem modelo fechado: o estudo inteiro
roda com ferramentas abertas e modelos de pesos abertos servidos
gratuitamente, sobre uma base comercial real.

O contexto semântico atua como instrução, e não como garantia. Ele
aumenta a proporção de acertos, mas o modelo pode deixar de aplicar uma
regra presente, e um contexto mal especificado, como um sinônimo
genérico, pode piorar o resultado em vez de melhorá-lo. O efeito de uma camada
semântica depende, portanto, da precisão com que ela é escrita e da
consistência das consultas usadas para avaliá-la. Seu custo é um prompt
maior em cada chamada, que aumenta a latência e o consumo de tokens.

Os objetivos específicos foram atingidos: o protocolo pareado e
reprodutível foi construído sobre uma base real descaracterizada; as três
condições foram comparadas em dois modelos de famílias distintas, com
testes pareados e ajuste para comparações múltiplas; os erros foram
classificados e associados ao nível de contexto que os resolve; o custo
em latência foi medido; e o contexto semântico foi relacionado às
ferramentas industriais e ao benchmark da comunidade Databricks.

% -----------------------------------------------------------------------------
\section{Trabalhos futuros}
\label{sec:trabalhos-futuros}

Os resultados apontam quatro passos que decorrem diretamente das
limitações e dos achados deste estudo.

O primeiro é repetir o experimento com um banco de perguntas maior,
escrito por outras pessoas depois do congelamento do contexto. Um banco
maior aumenta o poder dos testes pareados, necessário para decidir se o
efeito existe no \arq{gpt-oss-120b}, e um banco escrito por terceiros
afasta a possibilidade de o vocabulário das regras ter sido ajustado ao
das perguntas. A mesma revisão deve tratar as perguntas cujo enunciado
admite mais de uma leitura (Q019, Q035 e Q040) e as convenções de tipo
temporal da classe D, que hoje limitam a acurácia máxima observável.

O segundo é comparar o contexto semântico textual com uma camada
semântica consultada por meio de um compilador, como as do dbt ou do
Cube. Os erros de referência órfã e de colunas extras, que mais
contribuem para o ganho, são evitados por construção nesse tipo de
arquitetura; a comparação mostraria quanto do ganho depende de o modelo
escrever o SQL.

O terceiro é estender o estudo a esquemas maiores, em que o contexto
inteiro não cabe no prompt. Nesse cenário, é preciso selecionar as
tabelas e as regras relevantes para cada pergunta, e a seleção passa a
ser parte do que se avalia, com efeito também sobre o custo em tokens e
a latência.

O quarto é incluir modelos de outros tamanhos e famílias, inclusive
modelos proprietários e modelos com menos de 10 bilhões de parâmetros,
para separar o efeito do tamanho do efeito da família, o que os dois
modelos avaliados aqui não permitem.
